% Part: computability
% Chapter: recursive-functions
% Section: other-recursions

\documentclass[../../../include/open-logic-section]{subfiles}

\begin{document}

\olfileid{cmp}{rec}{ore}
\olsection{Wangun Rekursi Liyane}

Kanthi nggunakake pasangan lan rerangken, kita bisa
mbenerake wangun rekursi primitif kang luwih ora lumrah
(lan migunani). Umpamane, kerep migunani yen rong
fungsi ditegesi bebarengan, kaya ing definisi iki:
\begin{align*}
h_0(\vec x, 0) & = f_0(\vec x) \\
h_1(\vec x, 0) & = f_1(\vec x) \\
h_0(\vec x, y+1) & = g_0(\vec x, y, h_0(\vec x, y), h_1(\vec x, y)) \\
h_1(\vec x, y+1) & = g_1(\vec x, y, h_0(\vec x, y), h_1(\vec x, y))
\end{align*}
Iki tuladha \emph{rekursi bebarengan}. Cara liya
kang migunani kanggo netepake fungsi yaiku menehi
nilai $h(\vec x, y+1)$ adhedhasar \emph{kabeh}
nilai $h(\vec x, 0)$, \dots,~$h(\vec x, y)$,
kaya ing definisi iki:
\begin{align*}
h(\vec x, 0) & = f(\vec x) \\
h(\vec x, y+1) & = g(\vec x, y, \tuple{h(\vec x, 0), \dots, h(\vec x, y)}).
\end{align*}
Pola iki nyatakake gagasan mau kanthi luwih ringkes:
\[
h(\vec x, y) = g(\vec x, y, \tuple{h(\vec x, 0), \dots, h(\vec x, y-1)})
\]
kanthi pangerten yen argumen pungkasan kanggo $g$
iku rerangken kosong nalika $y$ iku $0$.
Ing kaloro wangun kasebut, gagasane yaiku nalika
ngitung 'langkah penerus,' fungsi $h$ bisa
nggunakake kabeh rerangken nilai kang wis
diitung sadurunge. Iki diarani \emph{rekursi
adhedhasar kabeh nilai sadurunge}. Minangka
tuladha tartamtu, cara iki bisa mbenerake
definisi jinis iki:
\begin{align*}
h(\vec x, y) & = \begin{cases}
  g(\vec x, y, h(\vec x, k(\vec x, y))) & \text{yen $k(\vec x, y) < y$} \\
  f(\vec x) & \text{yen ora}
\end{cases}
\end{align*}
Kanthi tembung liya, nilai $h$ ing $y$ bisa
diitung adhedhasar nilai $h$ ing \emph{sembarang}
nilai sadurunge kang diwenehake dening~$k$.


\begin{prob}
  Netepna fungsi turah $r(x,y)$ kanthi rekursi
  adhedhasar kabeh nilai sadurunge. (Yen $x$, $y$
  wilangan asli lan $y > 0$, $r(x,y)$ iku wilangan
  kang luwih cilik tinimbang~$y$ kanthi syarat
  $x = z\times y + r(x,y)$ kanggo sawenehing~$z$.
  Supaya cetha, yen $y=0$, $r(x,0) = 0$.)
\end{prob}

Pikirna carane ngasilake fungsi-fungsi iki kanthi
rekursi primitif lumrah. Wangun pungkasan saka
rekursi primitif luwih luwes amarga
\emph{parameter} (nilai ing pinggir) bisa
diowahi saklumakune proses:
\begin{align*}
h(\vec x, 0) & = f(\vec x) \\
h(\vec x, y+1) & = g(\vec x, y, h(k(\vec x), y))
\end{align*}
Wangun iki uga bisa ditiru nganggo rekursi
primitif lumrah. (Nindakake iki ora gampang.
Minangka pituduh, coba urai itungane kanthi
tangan.)

\end{document}
